% Source: https://artofproblemsolving.com/community/c3h1491971_number_sequences
% Stable identity: aops-topic-1491971-q01
A decrease of $6\%$ multiplies the membership by $0.94$ each year. The starting number $300$ belongs to ``year $0$''; after $1$ year there are $300\cdot0.94=282$ members, after $2$ years $300\cdot0.94^2$, and so on. After $5$ years the formula gives
\[
300\cdot0.94^5=300\cdot0.7339040224=220.17\ldots,
\]
so the club had $\boxed{220}$ members.

In general, a quantity $P$ that changes by a fixed percentage $r$ each period equals $P(1\pm r)^t$ after $t$ periods; the initial value is not counted as one of the periods.
